\documentclass[11pt]{article}
\usepackage[T1]{fontenc}
\usepackage{amsmath,amssymb,amsthm}
\usepackage[margin=1in]{geometry}
\usepackage{graphicx}
\usepackage{microtype}
\usepackage{xcolor}
\usepackage[hidelinks]{hyperref}

\newtheorem{theorem}{Theorem}
\newtheorem{proposition}[theorem]{Proposition}
\newcommand{\Prob}{\operatorname{Prob}}

\title{Stationary probability kernels need not be disintegrations\\
\large A finite counterexample to Remark 3.13}
\author{Automated Math Discovery Run \texttt{fa\_banach\_001}, lane 8}
\date{11 August 2026}

\begin{document}
\maketitle

\begin{center}
\fcolorbox{black}{red!5}{\parbox{0.88\linewidth}{
\textbf{Result.} The answer to Amrutam--Kl\"otzer--Oppelmayer's
Remark~3.13 is negative.  Stationarity and the barycenter identity do not
force a measurable probability kernel to be the disintegration of a factor
map, even for finite spaces and trivial actions.
}}
\end{center}

\section{Source question}

For a separable $\Gamma$-$C^*$-algebra $\mathcal A$ and a measurable
$\Gamma$-space $(Y,\eta)$, Proposition~3.12 of \cite{AKO} produces a
$\mu$-stationary measurable field of states.  Taking
$\mathcal A=C(X)$ gives a measurable map
\[
 \varphi:Y\longrightarrow\Prob(X),
\]
and its barycenter is
\[
 \nu=\int_Y\varphi(y)\,d\eta(y).
\]
Remark~3.13 asks whether $\varphi$ is the disintegration map associated with
some factor map $\pi:(X,\nu)\to(Y,\eta)$.

\begin{figure}[ht]
\centering
\includegraphics[width=0.98\linewidth]{figures/open_problem_crop.png}
\caption{Proposition 3.12 and Remark 3.13 in arXiv:2405.17122, PDF page 12.}
\end{figure}

\section{Idea of the counterexample}

A disintegration has a support condition absent from stationarity: if
$\{\nu_y\}$ disintegrates $\nu$ along $\pi$, then
\[
 \nu_y(\pi^{-1}(y))=1
 \quad\text{for $\eta$-almost every }y.
\]
With trivial group actions every probability kernel is stationary.  We may
therefore choose a kernel whose measures overlap completely across distinct
values of $y$, making fiber support impossible.  A singleton domain already
rules out the existence of any factor; a two-point variant shows that the
failure persists even when factor maps do exist.

\section{Counterexamples}

\begin{theorem}\label{thm:main}
There are finite compact $\Gamma$-spaces $X,Y$, an admissible non-degenerate
probability measure $\mu$ on $\Gamma$, and a $\mu$-stationary measurable map
$\varphi:Y\to\Prob(X)$ such that, for
$\nu=\int_Y\varphi(y)\,d\eta(y)$, no factor map
$\pi:(X,\nu)\to(Y,\eta)$ has $\varphi$ as its disintegration.  One may
choose the example so that no factor exists at all, or so that factors exist
but the prescribed field is not their disintegration.
\end{theorem}

\begin{proof}
Let $\Gamma=\mathbb Z/2\mathbb Z$ and
$\mu=\tfrac12(\delta_e+\delta_g)$.  This measure is admissible and
non-degenerate: it is absolutely continuous with respect to Haar measure and
its support generates $\Gamma$.  Give every space below the trivial
$\Gamma$-action.

For the first example, let $X=\{*\}$ with its unique probability measure and
let $Y=\{0,1\}$ with
$\eta=\tfrac12(\delta_0+\delta_1)$.  There is only one map
$\varphi:Y\to\Prob(X)$, namely
\[
 \varphi(0)=\varphi(1)=\delta_*.
\]
It is $\mu$-stationary because both actions are trivial.  Its barycenter is
$\nu=\delta_*$.  Any measurable map $\pi:X\to Y$ is constant, so
$\pi_*\nu$ is a Dirac mass and cannot equal $\eta$.  Thus there is no factor
map, and in particular $\varphi$ is not the disintegration of one.

For a version in which factors do exist, take
$X=Y=\{0,1\}$ and $\nu=\eta=\tfrac12(\delta_0+\delta_1)$.  Define the
constant field
\[
 \varphi(0)=\varphi(1)=\nu.
\]
Again it is stationary, and its barycenter is $\nu$.  A map
$\pi:(X,\nu)\to(Y,\eta)$ is measure preserving only if it is a bijection.
Its two fibers are singletons, and uniqueness of finite disintegration gives
\[
 \nu_y=\delta_{\pi^{-1}(y)}.
\]
But $\varphi(y)=\nu$ gives each singleton mass $1/2$, so it is not this
disintegration.  This proves both assertions.
\end{proof}

\section{Exact obstruction and scope}

The second example can be checked without invoking uniqueness: a conditional
measure over $y$ must assign mass one to $\pi^{-1}(y)$, while the constant
uniform field assigns that singleton mass $1/2$.  Thus the missing ingredient
is exactly fiber support (or, equivalently in this finite setting, the
idempotent/conditional-expectation structure of the Markov kernel).

The singleton example refutes the existence claim using only the hypotheses
of Proposition~3.12.  The equal-cardinality example additionally shows that,
even when an ambient factor is available, an arbitrary stationary fixed point
need not be its disintegration.  The result does not rule out positive
theorems under added compatibility or fiber-support hypotheses, nor does it
classify when one can select a different stationary fixed point that is a
disintegration.

\section{Verification and novelty notes}

All stationarity equations reduce to identities because the actions are
trivial.  All measures and pushforwards are displayed explicitly, and the
disintegration calculation is finite.  No computational experiment is used.

A bounded official-arXiv search through 11 August 2026 used arXiv id
2405.17122, the exact title, and combinations of ``stationary measurable
map'', ``disintegration'', and ``factor map''.  It found the source paper but
no later answer to Remark~3.13.  Expert review should verify that the literal
scope of the remark is the hypotheses of Proposition~3.12; the final scope
paragraph separates this from possible strengthened formulations.

\begin{thebibliography}{9}
\bibitem{AKO}
T.~Amrutam, M.~Kl\"otzer, and H.~Oppelmayer,
\emph{Relative stationary dynamical systems}, arXiv:2405.17122.
\end{thebibliography}

\end{document}
